\section{Discussion and conclusions}
The observed vehicle positions illustrate why model difference order, reference mean and extrapolation method must be kept distinct. For a position series, $d=3$ has a direct relation to the third temporal derivative---jerk---but neither the raw third difference nor its smoothed counterpart is an independent physical reference signal. A smooth-looking jerk series can still underrepresent actual sudden driver maneuvers.

The normalized coordinate $S\in[0,1]$ solves an important interpretability problem: it has consistent meaning across difference orders, regardless of the scale of the internal penalty $\lambda$, and its inverse map does not depend on the fixed mean of differences. The user-facing selection grid is always expressed in $S$, not an arbitrary independently sampled log-$\lambda$ coordinate. Pretranslation and reuse are computational optimizations, not new statistical models.

The principal practical conclusion must be qualified. A favorable GCV-selected $d=3$ one-second held-out result in this restricted sample is a descriptive illustration and cannot establish generalized vehicle-prediction superiority. The fifteen training/holdout origins overlap within five recorded segments and do not justify naive independent-window inference. Selection criteria may yield nearly identical smooth trends, and the observed position itself includes unmeasured camera noise.

The training-data mean-difference policy uses no future information, but CV scores computed conditionally on that plug-in mean are not exact re-estimated LOOCV. Profiled Gaussian marginal likelihood requires variance assumptions; it is not point-estimator MLE for $\tau$ or joint ML estimation of $\mu$. Restricted residual BIC remains an exploratory constrained criterion because the original unrestricted interpolation boundary is mathematically pathological.

Our extrapolation is an explicit local polynomial applied to the historical PLS fit. It is \emph{not} a predictive distribution, and it cannot anticipate unobserved changes in acceleration, traffic interactions or lane changes. Future research requires independent maneuvers, substantially more vehicles, genuine measurement-error characterization, and a dedicated dynamic/state-space forecast comparison. The seven fitting backends are benchmarked for single-curve cost and numerical agreement; complete-selection runtime and independent multiprecision accuracy are addressed by companion studies.

\section*{Reproducibility and data rights}
The US DOT/FHWA NGSIM dataset (DOI: 10.21949/1504477) is credited as the sole source of observed positions; the dataset portal identifies a CC BY-SA 3.0 license. Source query and extraction rules, tracked vehicle identifiers, segment starting timestamps, immutable checksums, chronological forecast rows, solver diagnostics, source-generation scripts and editable PGF vector figures are preserved under the case-study directory. The original official-position subset has not been artificially contaminated or pre-smoothed. Figure-derived jerk is explicitly estimated rather than observed.

\section*{Declarations}
This is an illustrative research draft. Authorship, correspondence, funding, conflicts and any journal submission require coauthor review and agreement.
